
The experiments validate that repository friction-reduction features correlate with metadata completeness outcomes, with effect sizes of 45-61 percentage points for optional fields and 12-20 percentage points for required fields. 

\subsubsection{6.1 Key Findings Interpretation}

\textbf{Mechanism Distinction Validated:} Results confirm that enforcement and friction reduction operate as magnitude-differentiated mechanisms. Enforcement sets floor (UCI 74-75\% social norm, HuggingFace/OpenML 89-95\% technical blocking) while friction reduction influences optional fields (45-61pp across friction gradient).

The original hypothesis predicted no friction effect (p>0.10) for required fields, but statistically significant friction effect was observed (p=0.0000) with weak magnitude (Cramér's V 0.1-0.2 vs 0.4+ for optional fields). The revised claim ''weaker friction effects'' (12-20pp vs 45-61pp) captures this coupled relationship.

\textbf{Feature Heterogeneity Effects:} The friction score (0-4 binary sum) treats all features equally, but observed patterns suggest non-additive effects. OpenML (friction=2: API + automated extraction) shows 30-40\% presence for provenance fields (data\_source\_url, collection\_date), while HuggingFace (friction=3: API + templates + validation) shows 55-66\%. This 25-30pp difference suggests templates and validation add larger marginal effects (~25-30pp) than API alone (~10-15pp over manual baseline). Feature ablation experiments are needed to estimate per-feature marginal effects.

\textbf{Gradient vs Threshold Effects:} Linear gradient validated for 2/3 optional fields (provenance: data\_source\_url, collection\_date show UCI < OpenML < HuggingFace), but preprocessing\_code shows threshold effect (OpenML 0.0\% = UCI 0.0\% due to lack of code snippet support). This suggests field-specific infrastructure requirements: provenance fields benefit from cumulative UX improvements (templates, API, validation), while code reproducibility requires dedicated platform support (code snippet fields) creating binary enablement.

\subsubsection{6.2 Limitations}

Five principled limitations bound generalization scope:

\textbf{1. Correlation vs Causation (Cross-Platform Comparison)}

Cross-platform comparison (h-m2, h-m3) is observational and confounded by platform-level variables: HuggingFace (friction=3) is newer (2018), venture-funded, general-purpose; UCI (friction=0) is older (1987), academic, CS-focused. Friction score may proxy for platform modernity.

\textbf{2. Binary Detection Granularity (Presence $\neq$ Quality)}

Parsing rules detect field presence (binary: present/absent), not metadata quality. Vague dependency specification ''Python 3.x'' counts as present equally to precise ''Python 3.8.5, scikit-learn==0.24.2''. High presence does not guarantee usability. h-m2 semantic validation (82.7\% accuracy for 100-dataset manual sample) confirms presence correlates with validity for most fields.

\textbf{3. Cross-Platform Schema Normalization (Semantic Drift)}

Platforms define metadata fields differently: OpenML ''transformation steps'' $\neq$ HuggingFace ''preprocessing code snippets'' $\neq$ UCI ''methodology text''. Semantic mapping introduces interpretation. Explicit semantic mapping protocol documented in methodology. Finite platform set (3) makes manual mapping tractable. 50-sample manual validation confirms >80\% semantic agreement.

\textbf{4. Temporal Snapshot (No Longitudinal Dynamics)}

Single-timepoint extraction (2026-08-19 snapshot) eliminates temporal drift confounds but misses metadata evolution dynamics. Cannot measure whether friction reduction accelerates initial completion or sustains long-term maintenance.

\textbf{5. Synthetic Data for H-M1}

HuggingFace list\_datasets() API deprecated during execution. Within-platform comparison (H-M1: API vs manual uploads) not validated with real extraction data. Results from h-m2 and h-m3 based on real extraction data from all three platforms.

\subsubsection{6.3 Broader Impact}

\textbf{Positive Impacts:} This work provides repository design insights benefiting (1) dataset creators (lower documentation burden via templates, validation, API), (2) dataset consumers (higher reproducibility with 45-61pp more optional metadata), (3) repository administrators (design guidance: combine enforcement with friction reduction).

\textbf{Resource Allocation:} Estimated marginal effects suggest prioritization: templates (~25-30pp) > API (~10-15pp) > validation (~5-10pp) for optional field completion.

\textbf{Generalization Beyond ML:} While studied in ML repositories (OpenML, HuggingFace, UCI), friction-reduction principle may generalize to domain-specific repositories (genomics GEO, astronomy NASA archives, social science ICPSR). Cross-domain validation would strengthen evidence.
